% Integration-ready proof module; no canonical file has been edited.
% Requires amsmath, amssymb, amsthm, and hyperref in a standalone host.
% All coordinates and quadratic-form conventions are retained explicitly.

\section{The marked Niemeier lattice, its twelve theta indices, and the umbral shadow}
\label{nus:sec:shadow}

The original objects are
\[
 A=\{(x_0,\ldots,x_5)\in\mathbb Z^6:\sum_jx_j=0\},\qquad
 D=\{d\in\mathbb Z^4:\sum_id_i\equiv0\pmod2\},\qquad
 R=A^{\perp4}\perp D.
\]
The metrics are the displayed Euclidean metrics.  The four copies of $A$
retain their original fibre order $f_i289^j$ for $0\le j<6$, with fibre
labels $(00,10,01,11)$.  Their discriminant generator is
$\lambda=(5,-1,-1,-1,-1,-1)/6$.  The discriminant of $D$ retains
$10\mapsto v=(1,0,0,0)$, $01\mapsto s=(1,1,1,1)/2$, and
$11\mapsto c=(-1,1,1,1)/2$.  Consequently the original norm-valued
quadratic refinement, with codomain $\mathbb Q/2\mathbb Z$, is
\[
 q_R(k_1,k_2,k_3,k_4;y)
 =\frac56\sum_i k_i^2+\mathbf1_{y\ne0}\pmod{2\mathbb Z}.
\]
The original glue and Niemeier lattice are
\[
 \begin{split}
 E&=\{u\in\mathbb F_2^4:\sum_i u_i=0\},\qquad
 \tau(u)=(u_2+u_4,u_3+u_4),\\
 j(a,b)&=(a+b,a+2b,a,b),\qquad C_3=j(\mathbb F_3^2),\\
 H&=\{(3u+2j(a,b);\tau(u)):u\in E,\ (a,b)\in\mathbb F_3^2\},\\
 N&=\{x\in R^\vee:x+R\in H\}.
 \end{split}
\]
Thus $|H|=8\cdot9=72$ and $N/R=H$, with the root system
$X=A_5^4D_4$ proved in the retained lattice calculation.  In particular,
no norm-two vectors are introduced by the nonzero glue classes: their
original minima are four or six.  We use this proved literal root system,
rather than a rank-only identification, throughout the construction below.

\subsection{The primary prescription and its complete coefficient matrix}

Write $e(t)=\exp(2\pi\mathrm i t)$ and $q=e(\tau)$ for
$\tau\in\mathbb H$.  Cheng--Duncan--Harvey, arXiv:1307.5793v2,
Section~3.2 \cite[Sections~3.2--3.3]{esn:CDH2014}, source labels
\texttt{thetexp}, \texttt{transf\_theta},
and \texttt{transf\_theta\_2}, use
\[
 \theta_{m,r}(\tau,z)
 =\sum_{k\equiv r\ (2m)}q^{k^2/(4m)}e(kz),
 \qquad r\in\mathbb Z/(2m).
\]
Their Section~3.3, source label \texttt{def:OmegaMatrices}, defines
\[
 \Omega_m(n)_{r,s}
 =\begin{cases}
 1,&r+s\equiv0\pmod{2n},\quad r-s\equiv0\pmod{2m/n},\\
 0,&\text{otherwise},
 \end{cases}
 \qquad n\mid m.
\]
Their ADE table, source label \texttt{ADE1}, assigns
\[
 \Omega^{A_{m-1}}=\Omega_m(1),\qquad
 \Omega^{D_{m/2+1}}=\Omega_m(1)+\Omega_m(m/2),
\]
and extends by addition over components with the same Coxeter number.
These are the primary authors' prescribed maps.  The following calculation
specializes and proves their values in the present marked coordinates; it
does not assert a new ADE classification.

The six-cycle $e_j\mapsto e_{j+1\bmod6}$ on an $A$ fibre has
characteristic polynomial $(x^6-1)/(x-1)$ on its augmentation space.
Its order is six.  In $D$, take the original simple roots
$\alpha_1=e_1-e_2$, $\alpha_2=e_2-e_3$, $\alpha_3=e_3-e_4$,
$\alpha_4=e_3+e_4$.  Their reflection product
$s_{\alpha_1}s_{\alpha_2}s_{\alpha_3}s_{\alpha_4}$ sends
\[
 e_1\mapsto e_2,\quad e_2\mapsto e_3,\quad
 e_3\mapsto-e_1,\quad e_4\mapsto-e_4.
\]
It has order six and characteristic polynomial
\[
 (x^3+1)(x+1)
 =(x^2-x+1)(x+1)^2
 =\frac{(x^2-1)(x^6-1)}{(x-1)(x^3-1)}.
\]
Thus the original five components have Coxeter number $h=m=6$, and the
primary theta prescription has precisely $2h=12$ indices and denominator
$4h=24$ in its exponents.  The root rank remains $4\cdot5+4=24$;
the glue index remains $72$.

Put $J=\Omega_6(3)$.  Its two congruences are
\[
 r+s\equiv0\pmod6,\qquad r-s\equiv0\pmod4.
\]
The simultaneous solutions modulo twelve consist of the single value
$s\equiv5r\pmod{12}$: this value satisfies both congruences, and the
difference of any two solutions is divisible by both six and four, hence
by twelve.  Therefore
\[
 \boxed{\Omega^X=5I_{12}+J,\qquad J_{r,s}=\mathbf1_{s\equiv5r\ (12)}.}
\]
In the full original index order $0,1,\ldots,11$ this is
\begingroup
\csname c@MaxMatrixCols\endcsname=12\relax
\[
\Omega^X=\begin{pmatrix}
6&0&0&0&0&0&0&0&0&0&0&0\\
0&5&0&0&0&1&0&0&0&0&0&0\\
0&0&5&0&0&0&0&0&0&0&1&0\\
0&0&0&6&0&0&0&0&0&0&0&0\\
0&0&0&0&5&0&0&0&1&0&0&0\\
0&1&0&0&0&5&0&0&0&0&0&0\\
0&0&0&0&0&0&6&0&0&0&0&0\\
0&0&0&0&0&0&0&5&0&0&0&1\\
0&0&0&0&1&0&0&0&5&0&0&0\\
0&0&0&0&0&0&0&0&0&6&0&0\\
0&0&1&0&0&0&0&0&0&0&5&0\\
0&0&0&0&0&0&0&1&0&0&0&5
\end{pmatrix}.
\]
\endgroup
The unary series and the actual shadow are
\[
 S_{6,r}(\tau)=\frac1{2\pi\mathrm i}
  \left.\frac{\partial\theta_{6,r}}{\partial z}\right|_{z=0}
 =\sum_{k\equiv r\ (12)}kq^{k^2/24},
 \qquad S^X=\Omega^X S_6.
\]
All series and their derivatives converge normally on compact subsets of
their stated domains, by Gaussian decay.  Substitution $k\mapsto-k$
proves $S_{6,-r}=-S_{6,r}$, so $S_{6,0}=S_{6,6}=0$.
Let $\mathcal A:\mathbb C^5\to\mathbb C^{12}$ insert zero in positions
zero and six and put $(\mathcal Au)_{12-r}=-u_r$ for $1\le r\le5$.
This retains, and explicitly reconstructs, all twelve components.  Then
\[
 \Omega^X\mathcal A=\mathcal A M^X,\qquad
 \boxed{M^X=\begin{pmatrix}
 5&0&0&0&1\\0&4&0&0&0\\0&0&6&0&0\\
 0&0&0&4&0\\1&0&0&0&5
 \end{pmatrix}.}
\]
Indeed its $(r,s)$ entry is
$\Omega^X_{r,s}-\Omega^X_{r,12-s}$; direct substitution gives the matrix.
In particular the fourth A-type multiplicity and the full D-type
contribution have both been retained.  The diagonal entries
$(5,4,6,4,5)$ are exactly the Coxeter exponent multiplicities:
four copies of $(1,2,3,4,5)$ and one copy of $(1,3,3,5)$.

\subsection{Quadratic refinements and the finite theta transformations}

Let $L=\mathbb Z\ell$ with $(\ell,\ell)=12$.  Its dual is
$L^\vee=\mathbb Z\ell/12$.  Hence
\[
 A_L=L^\vee/L=\mathbb Z/12,\qquad
 q_L(r)=r^2/12\pmod{2\mathbb Z},\qquad
 Q_L(r)=q_L(r)/2=r^2/24\pmod{\mathbb Z},
\]
and the associated bilinear form is $B_L(r,s)=rs/12\pmod{\mathbb Z}$.
The symbols $q_L$ and $Q_L$ deliberately have different codomains and
retain the factor two.

Define the unitary Fourier matrices and reflection by
\[
 (F_\pm)_{r,s}=12^{-1/2}e(\pm rs/12),\qquad
 R_{r,s}=\mathbf1_{s\equiv-r\ (12)}.
\]
Finite geometric-series summation gives
\[
 F_\pm^2=R,\qquad F_-=RF_+=F_+R,\qquad F_+F_-=I.
\]
The substitution $k\mapsto k+12$ shows that $e(k^2/24)$ depends only on
$r=k\bmod12$.  Thus, with
$D_{r,s}=e(r^2/24)\delta_{r,s}$,
\[
 \theta_6(\tau+1,z)=D\theta_6(\tau,z).
\]
Poisson summation of the Gaussian gives
\begin{align*}
 \theta_6(-1/\tau,z/\tau)
 &=\sqrt{-\mathrm i\tau}\,e(6z^2/\tau)F_-\theta_6(\tau,z),\\
 \theta_6(-1/\tau,-z/\tau)
 &=\sqrt{-\mathrm i\tau}\,e(6z^2/\tau)F_+\theta_6(\tau,z).
\end{align*}
Here $\sqrt{-\mathrm i\tau}$ is the holomorphic square root positive
when $\tau$ is positive imaginary.  For completeness, apply the Fourier
transform convention $\widehat f(u)=\int_{\mathbb R}f(x)e(-ux)\,dx$
to $f(x)=\exp(\pi\mathrm i\tau x^2/12+2\pi\mathrm i zx)$.
Completing the square evaluates its transform; the real positive Gaussian
integral is obtained by squaring the integral and using polar coordinates,
and holomorphic continuation gives the expression for $\Im\tau>0$.
Periodizing $f(r+12x)$, its $n$-th Fourier coefficient is
$12^{-1}e(nr/12)\widehat f(n/12)$.  Summing these coefficients at zero
and grouping $n$ by its residue modulo twelve gives the first displayed
transformation.  Absolute Gaussian convergence justifies the integration,
grouping, and Fourier evaluation.  Replacing $z$ by $-z$ and using
$\theta_{6,-r}(\tau,z)=\theta_{6,r}(\tau,-z)$ gives the second.
The second, with its plus-sign Fourier matrix, is precisely CDH's
convention; the first uses the usual generator
$S=\left(\begin{smallmatrix}0&-1\\1&0\end{smallmatrix}\right)$
with elliptic argument $z/\tau$.

The matrices carrying the finite multiplier for this first convention are
\[
 U(T)=D,\qquad U(S)=e(-1/8)F_-.
\]
The square-root weight factor has not been included in $U$.
Differentiation of the theta transformation at $z=0$ gives
\[
 S_6(-1/\tau)=\tau\sqrt{-\mathrm i\tau}F_-S_6(\tau),\qquad
 S_6(\tau+1)=D S_6(\tau).
\]
Equivalently, CDH's reflected convention gives
$S_6(-1/\tau)=-\tau\sqrt{-\mathrm i\tau}F_+S_6(\tau)$.
The equality follows from $RS_6=-S_6$.  Thus the unary series has weight
$3/2$ with the same finite multiplier; theta has weight $1/2$.

The permutation $J$ preserves $Q_L$ because
$((5r)^2-r^2)/24=r^2\in\mathbb Z$.  Since $5^{-1}=5$ modulo twelve,
\[
 (JF_\pm)_{r,s}=12^{-1/2}e(\pm5rs/12)
 =(F_\pm J)_{r,s}.
\]
Therefore $\Omega^X$ commutes with both finite multiplier matrices.
The shadow transformation follows by applying this proved intertwiner.
Its component series vanish at infinity since every nonzero summand has
strictly positive exponent; the modular transformation shows the same
cuspidal behavior at every cusp of the full modular group.

There is also an exact comparison with the original A5 discriminant,
without replacing its form.  Reduction
\[
 p:\mathbb Z/12\longrightarrow\mathbb Z/6,\qquad p(r)=r\bmod6
\]
has kernel $\{0,6\}$ and satisfies
\[
 q_A(p(r))=\frac56r^2=10q_L(r)\pmod{2\mathbb Z}.
\]
Thus it is a surjective quadratic similitude with precisely multiplier
ten, and its bilinear forms satisfy the same multiplier relation.
It is not asserted to be an isometry.  Applying $p$ to four coordinates
and the identity to the retained $D$ discriminant gives
\[
 \widetilde D=(\mathbb Z/12)^4\oplus\mathbb F_2^2
 \longrightarrow D_R,
 \qquad (r;y)\longmapsto(p(r_1),\ldots,p(r_4);y).
\]
The pullback quadratic refinement is exactly
$10\sum_iq_L(r_i)+\mathbf1_{y\ne0}$, its radical is
$\{0,6\}^4\oplus0$, and its quotient by that radical is the original
$(D_R,q_R)$.  To prove the radical claim, pairing with each coordinate
generator forces $5r_i/6\in\mathbb Z$, hence $6\mid r_i$; the
nondegenerate $D$ form forces $y=0$.  The vectors so obtained have
quadratic value $30$ in each nonzero coordinate, hence zero modulo two.
The inverse image of the original $H$ has order $16\cdot72=1152$ and
is isotropic for this pullback form.  This gives the exact relation of the
theta indices to all original glue coordinates, with the form and
multiplier specified on both sides.

\subsection{The full Fourier determinant and the retained order-twelve character}
\label{nus:subsec:determinant}

Let $\chi:\mathrm{SL}_2(\mathbb Z)\to\mathbb Z/12$ be the retained
marking
\[
 \chi(T)=-1,\qquad \chi(S)=3.
\]
These values respect the presentation
$S^4=I$, $(ST)^3=S^2$: in the abelian target the required relations
are $4\cdot3=0$ and $3(3-1)=2\cdot3$ modulo twelve.
They therefore define a character, and its value at $T$ makes it onto.
We now calculate the determinant of the entire twelve-dimensional
finite theta representation, rather than that of a reduced component
list.

Put $\zeta=e(1/12)$.  The Vandermonde determinant gives
\[
 \det F_+=12^{-6}\prod_{0\le j<k\le11}(\zeta^k-\zeta^j).
\]
For every such pair,
\[
 \zeta^k-\zeta^j
 =2\mathrm i\sin\!\left(\frac{\pi(k-j)}{12}\right)
       \exp\!\left(\frac{\pi\mathrm i(j+k)}{12}\right),
\]
whose sine is strictly positive.  There are $\binom{12}{2}=66$ pairs,
and each index occurs in eleven of them, so
\[
 \sum_{j<k}(j+k)=11\sum_{j=0}^{11}j=11\cdot66=726.
\]
The phase of the determinant is consequently
\[
 \mathrm i^{66}\exp(726\pi\mathrm i/12)=(-1)\mathrm i=-\mathrm i.
\]
Its remaining factor is positive.  Unitarity, proved by the finite
geometric sum above, makes its absolute value one.  It follows exactly
that
\[
 \boxed{\det F_+=-\mathrm i,\qquad\det F_-=\mathrm i.}
\]
In particular retaining CDH's reflection changes the Fourier determinant
by $\det R=(-1)^5=-1$; it cannot be dropped.

The sum of the squares of all twelve indices is
\[
 \sum_{r=0}^{11}r^2=\frac{11\cdot12\cdot23}{6}=506.
\]
Consequently
\[
 \boxed{\det U(T)=e(506/24)=e(1/12),\qquad
 \det U(S)=e(-12/8)\det F_-=-\mathrm i=e(-3/12).}
\]
For an explicit check that the finite matrices have the relevant
metaplectic relations, let
$G=\sum_{r=0}^{11}e(r^2/24)$.  Listing the squares modulo twenty-four
gives
\[
 G=1+4e(1/24)+2e(1/6)+2e(3/8)+2e(2/3)+e(1/2).
\]
The first and last terms cancel, as do the terms with arguments $1/6$
and $2/3$.  The elementary sine and cosine values at $\pi/12$ and
$3\pi/4$ give
\[
 4e(1/24)+2e(3/8)=\sqrt6(1+\mathrm i)
 =\sqrt{12}\,e(1/8).
\]
Completing the square in the finite sum proves
\[
 (F_-D)^2=e(1/8)D^{-1}F_-,\qquad
 (F_-D)^3=e(1/8)R.
\]
For the first equality, its $(r,s)$ entry before evaluation is
$12^{-1}e(s^2/24)\sum_te(t^2/24-t(r+s)/12)$; shifting $t$ by
$r+s$ produces exactly the stated Gauss sum and matrix entry.
The second follows by multiplication and $F_-^2=R$.
It follows that
\[
 (U(S)U(T))^3=U(S)^2=-\mathrm iR,\qquad U(S)^4=-I_{12}.
\]
Taking determinants kills this metaplectic central sign because the
dimension twelve is even.  The determinant therefore descends to an
ordinary character of $\mathrm{SL}_2(\mathbb Z)$.  Its two computed
generator values, together with generation by $S,T$, prove
\[
 \boxed{\det U(\gamma)=e(-\chi(\gamma)/12).}
\]
The dual determinant is its inverse and is exactly
$e(\chi(\gamma)/12)$, with the retained signs $T\mapsto e(-1/12)$,
$S\mapsto\mathrm i$.

This is a statement about the determinant multiplier, with its domain
and codomain specified.  It is not a statement that a vector of weight
$1/2$ has scalar weight one after taking a determinant.  The separate
Dedekind-eta identities
\[
 \eta(\tau+1)=e(1/24)\eta(\tau),\qquad
 \eta(-1/\tau)=\sqrt{-\mathrm i\tau}\,\eta(\tau)
\]
show that $\eta^2$ has weight one and multiplier $\det U$, whereas
$\eta^{-2}$ has weight minus one and multiplier $(\det U)^{-1}$.
Indeed $\eta^2(S\tau)=(-\mathrm i)\tau\eta^2(\tau)$ and
$\eta^{-2}(S\tau)=\mathrm i\tau^{-1}\eta^{-2}(\tau)$.
The integer-weight factors $(c\tau+d)^{\pm1}$ obey the cocycle identity
by matrix multiplication, so the two generator formulas prove the
statements for every $\gamma$.  At $-I$ the character and odd integer
weight each contribute a minus sign, whose product is one.

\subsection{The residue Fourier map and what the Weyl quotient retains}

The original marking is
$\kappa(289^j)=j\pmod6$, with the full ordered residue list
\[
 (289^j)_{j=0}^5=(1,289,361,169,121,529)\pmod{840}.
\]
On the six-dimensional complex coordinate carrier define
\[
 E_2:\mathbb C[S_{840}]\longrightarrow\mathbb C[\mathbb Z/12],
 \qquad E_2(e_{289^j})=e_{2j},
\]
and on a separately indexed $\mathbb C^6$ define
\[
 B(e_k)=\frac{e_k+e_{k+6}}{\sqrt2},\qquad0\le k<6.
\]
Both are injective and preserve the indicated coordinate Hermitian
products.  Retain the original negative-sign normalized Fourier matrix
$F_6^-=(6^{-1/2}e(-kj/6))_{k,j}$.  Entry-by-entry calculation gives
\[
 \boxed{F_-E_2=B F_6^-.}
\]
Indeed the $(r,j)$ entry of the left side is
$12^{-1/2}e(-rj/6)$, and the right side has the same value, using
$r\bmod6$.  Thus the original unnormalized sample transform is carried
into this formula by its exact factor $\sqrt6$.
Translation $j\mapsto j+1$ intertwines under $E_2$ with translation
$r\mapsto r+2$ modulo twelve.  This gives a literal Fourier and
translation map, not an identification of different quadratic forms.
In particular $Q_L(6)=3/2=1/2\pmod{\mathbb Z}$, so the quotient
$\mathbb Z/12\to\mathbb Z/6$ does not carry $Q_L$ itself to a
quadratic form on the quotient.  Its exact multiplier-ten comparison
with the original A5 form was proved above.

For the original lattice action, $\rho(289)$ is the same six-cycle in
each of the four $A$ coordinate sets and is the identity on $D$.
Each coordinate transposition of $A$ is reflection in the corresponding
root $e_i-e_j$, so the full residue action is in
$W(R)=W(A)^4\times W(D)$.  It acts trivially on every A discriminant
because
$Q_\sigma\lambda-\lambda=e_{\sigma(0)}-e_0\in A$, and it fixes
the D discriminant.  Consequently
\[
 S_{840}\xrightarrow{\rho}O(N)\longrightarrow O(N)/W(R)
\]
is the trivial homomorphism.  This proves exactly which residue-group
action is killed, without discarding its already constructed Fourier map.

The original marked quotient is
\[
 G=O(N)/W(R)\cong\mathrm{GL}_2(\mathbb F_3).
\]
Here its actual coordinate map uses
$\ell_1=(1,1)$, $\ell_2=(1,2)$, $\ell_3=(1,0)$,
$\ell_4=(0,1)$.  For $A\in\mathrm{GL}_2(\mathbb F_3)$ determine
the unique signs and permutation by
\[
 \ell_iA=\varepsilon_i\ell_{p(i)},\qquad
 (P_pu)_i=u_{p(i)},\qquad B_p\tau(u)=\tau(P_pu).
\]
The four row lines exhaust the projective line, so the signs and
permutation exist uniquely.  The binary equation defines $B_p$ uniquely
since $\tau$ is onto.  Then $(k;y)\mapsto(\varepsilon_i k_{p(i)};B_py)$
preserves precisely the original two glue codes, and its ternary
restriction is $M_gj=jA$.  Conversely a marked glue stabilizer gives this
same $A$.  Multiplication of these identities proves the group
isomorphism, with the original row/column convention retained.

To compare its chosen signed lifts with CDH's simple-root permutation
section, let $Q_{\rm rev}e_j=e_{5-j}$ in an $A$ fibre.  When
$\varepsilon_i=-1$, replace the representative $-I$ by
$-Q_{\rm rev}$; their quotient is $Q_{\rm rev}\in W(A)$ and
$-Q_{\rm rev}(e_j-e_{j+1})=e_{4-j}-e_{5-j}$.
Thus the resulting map permutes the original simple roots by diagram
reversal.  When $\varepsilon_i=1$ no change is needed.
For D, put $\Delta=(\alpha_1\ \alpha_2\ \alpha_3\ \alpha_4)$.
Associate the original classes $v,s,c$ with outer nodes $1,4,3$,
respectively; their fundamental weights are
$e_1$, $(1,1,1,1)/2$, and $(1,1,1,-1)/2$, the last of which differs
from the retained $c$ by $(1,0,0,-1)\in D$.
For $B_p$, let $P_{B_p}$ permute these outer nodes exactly as $B_p$
permutes $v,s,c$ and fix node two.  Then
\[
 d_{B_p}=\Delta P_{B_p}\Delta^{-1}
\]
is the exact original-coordinate diagram representative.  Its
orthogonality follows because permuting the outer nodes preserves the
original Gram matrix.  It preserves $D$, and its action on fundamental
weights proves its discriminant action is exactly $B_p$.
Its quotient with the originally chosen D lift has trivial discriminant
action, hence belongs to the proved Weyl kernel $W(D)$.
This specifies the actual change of section on every summand.

Define the following four class functions using those retained maps:
\[
 a_g=\#\{i:p(i)=i\},\quad
 b_g=\sum_{p(i)=i}\varepsilon_i,\quad
 s_g=\operatorname{sgn}(B_p|_{\{v,s,c\}}),\quad
 c_g=\#\operatorname{Fix}(B_p|_{\{v,s,c\}})-1.
\]
The first two are the traces on the four-component permutation and
signed permutation representations.  The last is the trace on the
two-dimensional sum-zero subspace of the three outer-node permutation
representation, and $s_g$ is its sign character.  Thus these are precisely
CDH's $(\bar\chi^{X_A},\chi^{X_A},\chi^{X_D},\check\chi^{X_D})$;
the remaining $\bar\chi^{X_D}$ is one because the single D component
is fixed.  The comparison of sections above proves this identification
for the actual marked quotient.

CDH Section~5.1 \cite[Section~5.1]{esn:CDH2014} defines the full
twisted coefficient matrix by
\[
 (\Omega_g^{X_A})_{r,s}
 =\bigl(b_g\mathbf1_{r\ \text{even}}+
          a_g\mathbf1_{r\ \text{odd}}\bigr)\delta_{r,s},
\]
and, for its exceptional $m=6$ D4 case, by
\[
 (\Omega_g^{X_D})_{r,s}=
 \begin{cases}
 c_g\delta_{r,s},&r\equiv0,3\pmod6,\\
 \delta_{r,s}+s_gJ_{r,s},&r\equiv1,5\pmod6,\\
 s_g\delta_{r,s}+s_gJ_{r,s},&r\equiv2,4\pmod6.
 \end{cases}
\]
All deltas in this formula have indices modulo twelve.  The exact maps
from the marked glue quotient to the prescribed twisted shadows are
\[
 g\longmapsto\Omega_g^X=\Omega_g^{X_A}+\Omega_g^{X_D},\qquad
 S_g^X=\Omega_g^XS_6.
\]
They are class-function constructions from traces of representations;
the assignment $g\mapsto\Omega_g^X$ is not asserted to be a group
representation.  On the odd five-component presentation the explicit
matrix is
\[
 \boxed{M_g^X=\begin{pmatrix}
 a_g+1&0&0&0&s_g\\
 0&b_g&0&0&0\\
 0&0&a_g+c_g&0&0\\
 0&0&0&b_g&0\\
 s_g&0&0&0&a_g+1
 \end{pmatrix}.}
\]
This follows by the same signed component reconstruction $\mathcal A$:
in rows two and four the two D terms cancel, in row three the contribution
is $c_g$, and rows one and five pair with one another.
For reference, all possible character tuples are
\[
\begin{array}{c|rrrr|r}
 [g]&a_g&b_g&s_g&c_g&\text{number of elements}\\\hline
 1A&4&4&1&2&1\\
 2A&4&-4&1&2&1\\
 2B&2&0&-1&0&12\\
 4A&0&0&1&2&6\\
 3A&1&1&1&-1&8\\
 6A&1&-1&1&-1&8\\
 8A,8B&0&0&-1&0&12
\end{array}
\]
with the CDH class names of their table
\texttt{tab:chars:eul:6}.  The formulas apply to each actual $A$, so
their validity does not depend on a choice of conjugacy-class names.
The finite checker obtains the table from all $48$ invertible matrices
over $\mathbb F_3$ and the original row forms.

The distinguished residue action has quotient $g=1$ and therefore
produces the untwisted shadow, while a product $\rho(r)L_g$ produces
the same $S_g^X$ for every $r$.  The four-component signed action and
the D4 triality quotient survive and enter the displayed matrix.  The
central element $-I\in\mathrm{GL}_2(\mathbb F_3)$, for example,
changes $b_g$ from four to minus four while fixing $a_g,s_g,c_g$;
it changes the even-index odd-space entries with exactly that sign.

The already marked composite
\[
 \mathrm{SL}_2(\mathbb Z)\xrightarrow{\chi}\mathbb Z/12
 \longrightarrow\mathbb Z/6\xrightarrow{j\mapsto289^j}S_{840}
\]
is onto, has kernel $\chi^{-1}(\{0,6\})$, and sends $T$ to
$289^5=529$.  Following it by the Weyl quotient kills it, by the
proved residue calculation.  Following $\chi$ into the dual determinant
multiplier retains its complete order twelve, by the determinant theorem.
These are explicit different outgoing morphisms from the same marking,
including their kernels and signs.

\subsection{The shadow-to-Fricke map with the original scalar seventy-two}

The Coxeter calculation above gives the original formal frame exponent
\[
 \pi_X=-5[1]+[2]-[3]+5[6].
\]
Let $\mathcal E(\sum a_d[d])=\prod_d\eta(d\tau)^{a_d}$.
It is a group homomorphism from addition of formal exponents to
multiplication of nonzero holomorphic functions, by the exponent laws.
The retained scalar function is
\[
 \boxed{T_6(\tau)=\mathcal E(-\pi_X)
 =\frac{\eta(\tau)^5\eta(3\tau)}
        {\eta(2\tau)\eta(6\tau)^5}.}
\]
The minus sign on $\pi_X$ is essential: $\mathcal E(\pi_X)=T_6^{-1}$.

Here is a direct coefficient map from the actual shadow matrix to this
eta function.  Its diagonal difference vector is
$\alpha=(5,4,6,4,5)$, and set $\alpha_0=0$ modulo six.
Define the Lambert-series map at this marked matrix by
\[
 \mathscr L(\Omega^X)(\tau)
 =1+\sum_{n\ge1}\alpha_{n\bmod6}
                    \frac{nq^n}{1-q^n}.
\]
Convergence is absolute and locally uniform since $|q|<1$ and
$\alpha$ is bounded.  The coefficient of $q^N$ is exactly
$\sum_{d\mid N}d\alpha_{d\bmod6}$; thus this specifies every
coefficient, not only its initial terms.
For all positive integers $n$ the retained multiplicities satisfy
\[
 \alpha_{n\bmod6}
 =5-\mathbf1_{2\mid n}+\mathbf1_{3\mid n}-5\mathbf1_{6\mid n}.
\]
Checking the six residues proves this identity.  The product definition
$\eta(\tau)=q^{1/24}\prod_{n\ge1}(1-q^n)$ and termwise logarithmic
differentiation then give
\[
 \boxed{f^X:=\mathscr L(\Omega^X)
 =-q\frac{d}{dq}\log T_6
 =5\lambda_6+\lambda_2-\lambda_3,}
\]
where the original functions on the right are
\[
 \lambda_d=q\frac{d}{dq}\log\frac{\eta(d\tau)}{\eta(\tau)}
 =\frac{d-1}{24}+\sum_{n\ge1}\sigma_1(n)(q^n-dq^{dn}).
\]
The constant is
$[5(6-1)+(2-1)-(3-1)]/24=1$, and the other coefficients agree by
the six-residue identity.  This proves the shadow-to-eta calculation
directly for every coefficient.

The precise primary connecting map is CDH
Section~4.2 \cite[Section~4.2]{esn:CDH2014}:
\[
 \sigma^X(\tau,z)=\sum_{r\ (12)}\overline{S_r^X(\tau)}
                                \theta_{6,r}(\tau,z),\qquad
 \mathscr S_{1,1}(\sigma^X)=f^X.
\]
Its source labels are \texttt{eqn:sigmaX}, \texttt{def:weight2},
\texttt{weight2:Aseries}, and \texttt{weight2\_EZ}.
The displayed Lambert map is its explicit specialization at this
root-system shadow.  It verifies the resulting function independently
of the general Shimura-lift theorem.  In particular the shadow and the
retained Fricke function are connected by this complete coefficient map.

The original modular subgroup here is $\Gamma_0(6)$, and
$[\mathrm{SL}_2(\mathbb Z):\Gamma_0(6)]=12$.  To compute the index,
reduction modulo six gives
$\mathrm{SL}_2(\mathbb Z/6)\cong
\mathrm{SL}_2(\mathbb F_2)\times\mathrm{SL}_2(\mathbb F_3)$ of
order $6\cdot24=144$.  Reduction is onto: elementary upper and lower
unipotent matrices generate over each field, and CRT choices of their
entries realize each factor while acting trivially on the other.
The upper triangular subgroup modulo six has $\varphi(6)\cdot6=12$
elements (choose its first unit diagonal entry and its upper-right entry;
the other diagonal entry is its inverse).  Its inverse image is
$\Gamma_0(6)$, proving index $144/12=12$.
This modular-group index and the twelve theta indices have their
specified constructions; neither replaces the lattice index seventy-two.

Under the exact Fricke transformation $W_6\tau=-1/(6\tau)$,
the eta inversion formula sends each $\eta(dW_6\tau)$ to
$\sqrt{-\mathrm i(6/d)\tau}\,\eta((6/d)\tau)$.
The retained exponent vector at $d=1,2,3,6$ is $(5,-1,1,-5)$;
its reversed vector is its negative and its exponent sum is zero.
The $\tau$ factors therefore cancel.  Since each $6/d$ is positive,
the consistent square-root branch gives the positive scalar
\[
 \prod_{d\mid6}(6/d)^{a_d/2}
 =6^{5/2}3^{-1/2}2^{1/2}=72.
\]
Consequently
\[
 \boxed{T_6(W_6\tau)=72/T_6(\tau).}
\]
For the weight-two slash action of the determinant-six representative,
\[
 (f|_2W_6)(\tau)=(6\tau^2)^{-1}f(-1/(6\tau)).
\]
Differentiating the last identity and using
$d(W_6\tau)/d\tau=(6\tau^2)^{-1}$ gives
\[
 \boxed{f^X|_2W_6=-f^X.}
\]
On the Jacobi side, the primary Eichler--Zagier operator
$\mathcal W_6(6)$ acts on theta by $\Omega_6(6)=R$.
This is verified by its defining congruences, $r+s=0\pmod{12}$,
$r-s=0\pmod2$.  On the odd shadow and on the corresponding
skew-Jacobi form its eigenvalue is minus one.  The specialized coefficient
map just proved carries that negative eigenline to the negative Fricke
eigenline of $f^X$.  This proves an exact equivariant linear comparison
on these lines; it does not evaluate an $\mathrm{SL}_2(\mathbb Z)$
character on the determinant-six matrix.

Finally, the original sheet coordinates remain
\[
 J_6=T_6+5+72/T_6,\qquad Y_6=T_6-72/T_6.
\]
Direct substitution and differentiation prove
\[
 J_6(W_6\tau)=J_6(\tau),\quad
 Y_6(W_6\tau)=-Y_6(\tau),\quad
 Y_6^2=(J_6-5)^2-288,\quad
 dJ_6=Y_6\,dT_6/T_6=-2\pi\mathrm iY_6 f^X\,d\tau.
\]
Thus the original scalar $72$, the additive shift $5$, the branch
constant $288$, and the differential factor $-2\pi\mathrm i$ all
survive the complete shadow-to-sheet calculation.

\subsection{Established scope and reproducibility}

The constructed maps are the ADE coefficient specialization, the exact
finite-theta intertwiner, the multiplier determinant and its dual, the
quadratic similitude to the original discriminants, the marked glue-group
trace prescription, the original six-to-twelve Fourier map, and the
coefficient-level shadow-to-Fricke map.  Their kernels, weights, signs,
and constants were specified and proved above.  They establish concrete
relations even where a direct isometry or an identical group action is
absent.  They do not construct a new graded moonshine module or a new
moonshine theorem.  The name ``umbral shadow'' here refers to precisely
the CDH prescription, with its source provenance retained.

The standard-library program \texttt{check\_shadow.py} uses exact
arithmetic in $\mathbb Q[z]/(z^8-z^4+1)$ with $z=e(1/24)$.
It verifies both full and odd matrices, the Fourier determinant by the
Vandermonde product, the Gauss sum, the character signs, the Fourier
embedding, all $48$ original-code character records, the first $240$
Lambert coefficients, and the squared positive Fricke scalar $72^2$.
The proofs above establish the unbounded identities independently of
those finite checks.
